\section{Introduction}
\label{sec:introduction}

Large Language Models (LLMs) have evolved from standalone text generation models into components of increasingly complex software systems. Modern LLM-based applications may incorporate prompt templates, retrieval-augmented generation (RAG), external tools, memory mechanisms, multi-agent coordination, iterative reasoning, and interaction with dynamic environments. Consequently, the quality of an LLM-based system can no longer be adequately characterized only by measuring the linguistic quality or factual correctness of its final response.

A failure observed at the output of an LLM system may originate at several different stages. An incorrect response may result from an ambiguous or incomplete user specification, inadequate retrieved context, an error in reasoning or tool selection, a failure during a multi-step agent trajectory, or limitations of the underlying model itself. Similarly, a response that appears correct according to an automated evaluator may still be incorrectly assessed if the evaluation mechanism is biased, unstable, poorly calibrated, or insufficiently aligned with human judgment. In production environments, additional factors such as latency, cost, availability, throughput, resource consumption, privacy, and safety may further determine whether a technically capable system can be considered satisfactory in practice.

This complexity has resulted in a rapidly expanding evaluation ecosystem. Numerous benchmarks and metrics have been proposed to assess individual properties of LLMs and LLM-based systems, while evaluation frameworks increasingly combine automated metrics, human feedback, LLM-based judges, execution-based tests, traces, and operational telemetry. At the same time, different research communities often employ distinct terminology for closely related concepts, and apparently similar metrics may operationalize different notions of quality.

A particularly visible example is the growing use of LLM-as-a-Judge approaches. LLM judges have become a common \emph{method} for evaluating model outputs, agent trajectories, and other artifacts. However, using an LLM as an evaluator is conceptually different from evaluating the quality of that evaluator. Reliability, consistency, bias, calibration, agreement with human experts, and sensitivity to evaluation prompts constitute a separate meta-evaluation problem. Conflating these two roles can obscure an important source of uncertainty in automated evaluation.

A similar distinction arises in end-to-end evaluation. A system may be evaluated from input to final task outcome without examining its internal stages. Such black-box evaluation is useful for measuring overall task success, but it provides limited information about the origin of failures. In contrast, other approaches explicitly relate quality across multiple stages, such as retrieval quality and response quality, or agent trajectory quality and task completion. A smaller subset attempts to combine several lifecycle stages together with evaluator reliability or operational characteristics. These approaches represent qualitatively different levels of evaluation integration.

This survey investigates this fragmented landscape from a system-oriented perspective. Rather than organizing the literature around individual benchmarks or particular evaluation technologies, we consider quality assessment across the lifecycle of an LLM-based system. Our central premise is that a comprehensive understanding of LLM quality requires distinguishing three complementary aspects: the object being evaluated, the quality dimension being measured, and the mechanism used to perform that measurement.

Accordingly, this work makes the following contributions:

\begin{itemize}
    \item We propose a lifecycle-oriented taxonomy of quality assessment for LLM-based systems, covering input and specification, context and retrieval, output, conversational interaction, agent processes and trajectories, and operational/system quality.
    \item We distinguish the \emph{quality dimension} being evaluated from the \emph{evaluation method} used to assess it, organizing methods such as deterministic metrics, semantic metrics, human evaluation, LLM-as-a-Judge, multi-judge evaluation, task-based evaluation, adversarial testing, trace-aware evaluation, and statistical reliability analysis.
    \item We explicitly model evaluator quality as a meta-evaluation problem, separating the use of an evaluator from evidence regarding its reliability.
    \item We introduce a taxonomy of evaluation integration that distinguishes isolated evaluation, cross-stage evaluation, black-box end-to-end evaluation, and lifecycle-integrated evaluation.
    \item We analyze how existing approaches combine quality dimensions, evaluation methods, and operational considerations, with the goal of identifying weakly connected areas despite the large number of existing metrics and frameworks.
\end{itemize}

The remainder of this article is organized as follows. Section~\ref{sec:background} introduces the conceptual foundations and terminology used throughout the survey. Section~\ref{sec:research-questions} presents the research questions. Section~\ref{sec:methodology} describes the literature search, eligibility criteria, classification protocol, and human validation procedure. Section~\ref{sec:taxonomy} presents the proposed taxonomy. The subsequent evidence-driven sections will synthesize evaluation methods, integration patterns, tools, and research gaps after completion of the corpus audit.
